\subsection{Isolated flight, turns, wind and battery}\label{sec:isolated}
\paragraph{Steady level flight and banked turns} At bank angle $\phi$ the load factor is $n_z=1/\cos\phi$ and the required lift coefficient $C_L=2mgn_z/(\rho\Va^2S)$; with the parabolic polar $C_D=C_{D0}+k_0C_L^2$ the drag, aerodynamic power $P_a=D\Va$ and electrical power $P_{\mathrm{el}}=P_a/\eta_p+P_{\mathrm{hotel}}$ follow. At $\Va=\SI{18}{m/s}$, $\phi=0$: $C_L=\nCLiso$, $C_D=0.0273$, $D=\SI{4.34}{N}$, $L/D=\nLoverD$, $P_a=\SI{78.0}{W}$ and $P_{\mathrm{el}}=\SI{\npIso}{W}$; induced drag is \ncdiShare\,\% of the total, which bounds what any wake can recover. At $\phi=25^\circ$ the load factor is 1.103, the turn radius $R_t=\Va^2/(g\tan\phi)=\SI{70.85}{m}$ and the load \SI{\npTurn}{W} (the preliminary version printed \SI{146.42}{W}, which its own polar does not give). The power curve over \SIrange{12}{26}{m/s} is shown in Fig.~\ref{fig:X6}: electrical power is minimal near \SI{\npowMinV}{m/s} and energy per kilometre is minimal near \SI{\npowRangeV}{m/s}, so the nominal \SI{18}{m/s} survey cruise is already close to the best-range speed.

\paragraph{Wind triangle} With wind vector magnitude $V_w$ blowing towards direction $\psi_w$ and a ground track $\chi$, the crab angle satisfies $\sin\Delta\psi=-V_w\sin(\psi_w-\chi)/\Va$ and the ground speed is $V_g=\Va\cos\Delta\psi+V_w\cos(\psi_w-\chi)$. Leg times are $L/V_g$ with airspeed (hence power) held constant. A headwind/tailwind pair therefore costs more energy than the zero-wind pair by the factor $[1-(V_{w\parallel}/\Va)^2]^{-1}$, which is evaluated leg by leg in the mission simulator (wind \SIrange{2}{6}{m/s}, direction uniform).

\paragraph{Turns between strips} Strips are $s=\SI{\ndyBase}{m}$ apart (65\,\% side overlap) or $s=\SI{74.0}{m}$ (the edge-optimised 14-strip plan). Because $s<2R_t=\SI{141.7}{m}$, a semicircle cannot connect adjacent strips; we compute the shortest curvature-bounded path \cite{dubins1957} (six candidate words LSL, RSR, LSR, RSL, RLR, LRL) from the end of one strip (heading north) to the start of the next (heading south) with a free straight run-in $y_e\ge0$ to the polygon edge, minimising total length; the result is \SI{441}{m} for $s=\SI{61.6}{m}$ and \SI{423}{m} for $s=\SI{74.0}{m}$, against \SI{222.6}{m} for the semicircle assumed in the preliminary version. The turn is flown at the nominal \SI{25}{\degree} bank.

\paragraph{Battery} The pack is six series cells of \SI{10}{Ah}; the open-circuit voltage is interpolated from a generic LiPo curve (\SI{3.30}{V} at 0\,\%, \SI{3.70}{V} at 20\,\%, \SI{3.86}{V} at 60\,\%, \SI{4.20}{V} at 100\,\% per cell) and the internal resistance is \SI{20}{m\ohm}. For a constant load $P$ the current solves $P=(V_{oc}-IR)I$, the cell-side power is $V_{oc}I\ge P$, and the state of charge decreases by $I\,\mathrm{d}t/Q$. Load-side energy $E_{\mathrm{load}}$ (the preliminary version's metric) and cell-side energy $E_{\mathrm{cell}}$ are both reported; the latter exceeds the former by about \SI{0.5}{\percent}. A \SI{20}{\percent} reserve policy gives a usable window of \SI{177.6}{Wh} (nominal \SI{22.2}{V}); the maximum per-aircraft mission energy is about \SI{38}{Wh} for the independent strategy, so the mission is far from the battery limit and the reserve never binds.

\subsection{Derivations behind the two corrections}\label{sec:derivations}
\paragraph{Root circulation} For an elliptic loading $\Gamma(y)=\Gamma_0\sqrt{1-(2y/b)^2}$ the lift is $L=\rho\Va\int\Gamma\,\mathrm{d}y=\rho\Va\Gamma_0\,\pi b/4$, and a pair of concentrated vortices of circulation $\Gamma_0$ spaced $b_0=\pi b/4$ apart carries the same lift and the same vertical impulse, which defines the roll-up of the elliptic wing. Hence
\begin{equation}\Gamma_0=\frac{L}{\rho\Va b_0}=\frac{\tfrac12\rho\Va^2SC_L}{\rho\Va\,\pi b/4}=\frac{2\Va SC_L}{\pi b}.\end{equation}
Dropping the $\tfrac12$ of the dynamic pressure gives $4\Va SC_L/(\pi b)$ (the preliminary Eq.), twice the correct value; for the present aircraft $\Gamma_0=\SI{\ngammaTrue}{m^2/s}$, not \SI{\ngammaPub}{m^2/s}. The vertical velocity of a vortex scales linearly with $\Gamma$, so the upwash and every first-order saving derived from it were overstated about twofold. A unit test compares the formula to $L/(\rho\Va b_0)$.

\paragraph{Rolling moment of the follower} A strip of the follower at spanwise station $y$ (measured from the follower centre line) sees an angle-of-attack increment $\Delta\alpha(y)=w(y)/\Va$; the section lift increment per unit span is $\Delta\ell'=q\,c(y)\,a\,\Delta\alpha$ with $q=\tfrac12\rho\Va^2$, the chord $c=S/b$ and the finite-wing lift slope $a=2\pi AR/(AR+2)$. The rolling moment about the longitudinal axis is $M_x=\int y\,\Delta\ell'\,\mathrm{d}y$ and its coefficient $C_l=M_x/(qSb)$, so
\begin{equation}C_{l,\mathrm{wake}}=\frac{a}{b^2}\int\frac{w(y)}{\Va}\,y\,\mathrm{d}y.\end{equation}
The preliminary version used $2\pi/(bS)$ in front of the integral, which has dimension $\mathrm{m}^{-1}$ and, for this aircraft, is \nclPubOverCorr$\times$ larger at $d_y=0.8b$; since the trim-drag penalty is proportional to $C_l^2$, the penalty zone and the knife-edge optimum of the preliminary version follow from the error.

\paragraph{From upwash to power} At matched airspeed and lift the induced-drag change is $\Delta C_{Di}=C_L\bar w/\Va$ (strip theory) or the change of the near-field bound-vortex drag $-\rho\sum\Gamma_j\Delta y_j w_j$ (vortex lattice); the net drag change subtracts the trim drag $k_{\mathrm{trim}}\delta_a^2$ with $\delta_a=-C_{l,\mathrm{wake}}/C_{l\delta_a}$. The follower's electrical power is obtained by inserting the new $C_D$ in the power chain above. Percentages in this paper always refer to the electrical power of an isolated aircraft at the same airspeed, so the hotel load dilutes the aerodynamic saving: a \SI{7.7}{\percent} drag reduction becomes \SI{6.2}{\percent} of electrical power.

\subsection{Planform vortex lattice for the CFD geometry}\label{sec:planform}
The CFD campaign uses an idealized wing-only geometry (Table~\ref{tab:cfdgeom}) rather than the rectangular-equivalent wing of the earlier models, so that CFD and vortex-lattice results are comparable on identical planforms (distinct from the cambered physical Skywalker-X8 airframe characterised in wind-tunnel and flight-identification studies \cite{gryte2018,reinhardt2022}). The planform vortex lattice places one horseshoe per span station with the bound vortex on the local quarter-chord line (swept polyline) and control points at $3c/4$; trailing legs are planar and parallel to the freestream. Forces use the Kutta--Joukowski relation; evaluating induced drag on the swept bound-vortex midpoints with a single chordwise panel gives $C_L=\nplanCL$ at $\alpha=3^\circ$, $C_{L\alpha}=\nplanCla$\,rad$^{-1}$ ($AR=\nplanAR$), and an apparent near-field Oswald factor of \nplanE{} ($>1$ due to single-panel swept bound-vortex self-upwash bias across the root kink, whereas far-field Trefftz-plane integration of the same trailing wake gives $e_{\mathrm{TP}}=\nplanETP$); near-field integration is retained only so that isolated and follower induced drag share identical bound-vortex quadrature, using differences rather than absolute drag. In the baseline cross-check, a single follower is solved one-way in the Biot--Savart velocity field of the isolated leader ($\Delta C_L=\nplanOneDCL$ at $d_x=3b$, $d_y=0.85b$). Because the OpenFOAM half-model enforces a symmetry plane at $y=0$ (Sec.~\ref{sec:cfdsetup}), we also solve the coupled symmetric three-wing VLM including the mirror-image second follower at $-d_y$ (centre separation $2d_y=1.7b=\SI{3.60}{m}$ at the nominal slot) and two-way upstream induction on the leader: mutual follower-to-follower upwash increases the follower lift increment to $\Delta C_L=\nplanSymDCL$ ($\nplanSymPct$\,\%), while upstream induction changes the leader's $C_L$ by only $\nplanSymLeadPct$\,\% (Sec.~\ref{sec:cfd}).
\begin{table}[H]\centering\small
\caption{Geometry of the CFD wing (clean wing, symmetric section, no twist, no winglets, no pod).}\label{tab:cfdgeom}
\resizebox{\textwidth}{!}{\begin{tabular}{@{}lll@{}}\toprule
Quantity & Value & Note\\\midrule
span $b$ & \SI{2.12}{m} & same as the aircraft model\\
root / tip chord & \SI{0.52}{m} / \SI{0.25}{m} & linear taper\\
leading-edge sweep & \SI{22}{\degree} & constant\\
planform area $S$ & \SI{\ncfdS}{m^2} & used as reference area in the CFD coefficients\\
mean aerodynamic chord & \SI{\ncfdMAC}{m} & $Re=\Va\,\overline{c}/\nu=\nReMAC\times10^{5}$\\
aspect ratio & \ncfdAR & $b^2/S$\\
section & NACA-0009-type thickness, blunt trailing edge at 98.5\,\% chord & thickness 9\,\% (root and tip)\\
reference point for moments & root quarter chord (leader), follower root quarter chord (follower) & \\\bottomrule
\end{tabular}}\end{table}

\subsection{OpenFOAM RANS set-up}\label{sec:cfdsetup}
\paragraph{Solver and models} Steady incompressible RANS with \texttt{simpleFoam} (OpenFOAM v1912, ESI; Ubuntu package 1912.200626) in SIMPLEC mode, $k$--$\omega$ SST turbulence closure, wall functions (\texttt{nutkWallFunction}, \texttt{omegaWallFunction}, \texttt{kqRWallFunction}), freestream $U_\infty=\SI{18}{m/s}$, $\nu=\SI{1.5e-5}{m^2/s}$, $\rho=\SI{1.225}{kg/m^3}$, turbulence intensity 0.1\,\% (inlet $k=1.5(IU)^2$, $\omega=\sqrt{k}/(C_\mu^{1/4}\ell)$ with $\ell=\SI{0.01}{m}$). Convection of momentum uses \texttt{linearUpwindV} with a cell-limited gradient, turbulence quantities use bounded upwind, the pressure equation GAMG (tolerance $10^{-8}$, relative $0.05$). Relaxation factors are 0.9 with the consistent (SIMPLEC) formulation. Convergence is monitored by residuals ($p$, $U$, $k$, $\omega$) and, for key cases, by the force coefficients at saved iterations.

\paragraph{Domain and boundary conditions} A half-model with a symmetry plane at $y=0$ through the leader's centre line represents the symmetric V formation (leader plus one follower stands for the pair of followers). The domain spans $x\in[-8,\,d_x+22]$\,m, $y\in[0,9]$\,m, $z\in[-7,7]$\,m (about $4b$ upstream, $\ge10b$ downstream of the follower, $4.2b$ lateral and $\pm3.3b$ vertical); inlet is fixed velocity, outlet fixed static pressure, and the lateral/vertical boundaries are slip walls, which is adequate for $\ge3b$ clearance and is checked indirectly by the mesh study. The leader wing is extended \SI{0.15}{m} through the symmetry plane so that the plane cuts the geometry cleanly; the follower is a full wing at $(d_x,d_y)$. Geometry rotates with $\alpha$ about the root quarter chord, the freestream stays along $+x$.

\paragraph{Meshing} Background hexahedral mesh of \SI{0.5}{m} cells; \texttt{snappyHexMesh} with surface refinement on the wings, distance-based refinement shells, a wake box along the leader's wake and a tip box around the wingtip vortex path (and around the follower). Three mesh levels (L0/L1/L2) change surface, wake and tip levels together (Table~\ref{tab:cfdmesh}); no prism layers are used, the first cell centre lies in the logarithmic region ($y^+\approx$ \nyplusZero, \nyplusOne and \nyplusTwo\ for L0, L1, L2), and the meshes pass \texttt{checkMesh} (maximum skewness $<2$, no negative volumes).

\paragraph{Force post-processing without function objects} The OpenFOAM v1912 binaries shipped with Ubuntu~24.04 abort in \texttt{functionObjectList} (an \texttt{OSHA1stream} error) whenever \emph{any} function object is present in \texttt{controlDict}, in serial and in parallel. Forces are therefore computed in Python from the exported patch fields (\texttt{foamToVTK}): the pressure force is $\rho\sum p\,\mathbf{n}\,A$ over the wing faces (outward normal of the fluid domain), the viscous force $\rho\sum(\nu+\nu_{t,w})\,\mathbf{U}_t/y\;A$ with the tangential owner-cell velocity $\mathbf U_t$ and wall distance $y$ of the owner cell centre---the same expression OpenFOAM evaluates as $\nu_{\mathrm{eff}}\,\partial\mathbf{U}/\partial n$ for wall-function patches. Leader loads are doubled for the half model. The routine is exercised by the unit tests and cross-checked on the isolated wing against the vortex-lattice lift (Sec.~\ref{sec:cfd}).
\begin{table}[H]\centering\small
\caption{OpenFOAM meshes (isolated half-wing, $\alpha=3^\circ$) and wall-function resolution.}\label{tab:cfdmesh}
\resizebox{\textwidth}{!}{\begin{tabular}{@{}lccccc@{}}\toprule
Level & surface levels (min,max) & wake box level & tip box level & cells (isolated) & $y^+$ mean (5--95\,\%)\\\midrule
L0 (coarse) & (4\,5) & 2 & 3 & 82\,728 & 87 (58--104) \\
L1 (medium) & (5\,6) & 3 & 4 & 381\,561 & 61 (42--73) \\
L2 (fine) & (6\,7) & 3 & 5 & 1\,251\,382 & 51 (30--87)
\\\bottomrule
\end{tabular}\end{table}
